\section{Additional Benefits of Proof Simplification}

\subsection{Training on Simplified Proofs Improves Generation} \label{subsec:training-simplified}

Next, we investigate whether fine-tuning on simplified proofs can be advantageous compared to fine-tuning on longer, raw proofs. To do so, we prepare two datasets of identical problems, the first containing a set of proofs generated by \texttt{Goedel-Prover-V2} and the second containing the same proofs simplified by ProofOptimizer-ExpIt. The average proof length of the original and simplified proofs is $147$ and $85$, respectively. We do continued supervised fine-tuning (SFT) starting from our base model (Sec.~\ref{subsec:lean-base-model}) with a standard negative log-likelihood (NLL) loss. 

In Fig.~\ref{fig:continued-training-plots} (left), we compare the training loss between the two datasets. As expected, the initial loss when using original proofs is higher, as models have not seen such long proofs during initial fine-tuning. However, the losses quickly converge. We observe that training on original proofs causes occasional loss spikes, which we suspect are due to several data batches that are hard to learn (e.g. extremely long proofs). Decreasing the learning rate mitigated these training loss spikes but did not improve validation accuracy. In Fig.~\ref{fig:continued-training-plots} (right), we compare the miniF2F scores of the two models during SFT, showing that training on simplified proofs results in slightly higher evaluation accuracy despite the two settings having identical training losses. 
%
\begin{figure}[H]
    \centering
    \begin{subfigure}{0.44\textwidth}
        \centering
        \includegraphics[width=\linewidth]{figures/continued-training/loss.pdf}
    \end{subfigure}
    \hfill
    \begin{subfigure}{0.44\textwidth}
        \centering
        \includegraphics[width=\linewidth]{figures/continued-training/minif2f.pdf}
    \end{subfigure}
    \caption{Training loss (left) and miniF2F score (right) after SFT on simplified vs. original proofs.}
    \label{fig:continued-training-plots}
\end{figure}
%
\subsection{Simplified Proofs Have a Shorter Execution Time}\label{subsec:proof-execution-time}
We also observe that proofs simplified by ProofOptimizer often exhibit a faster execution time. We measure proof execution time with \texttt{lake env lean --profile}, excluded library import time (imports are always the same but actual time may vary due to caching effects). We compare the execution times of each proof before and after iterative shortening in Fig. \ref{fig:speedup-plots} (scatter). For both datasets, we visibly observe that a majority of points lie below the $y=x$ line, signifying speedup. Fig. \ref{fig:speedup-plots} (histograms) also show the distribution of speedup ratios $\frac{\text{time}_\text{orig}}{\text{time}_\text{new}}$. Of the $75$ PutnamBench proofs, $50/75$ have a speedup of over $10\%$, and $22/75$ of those have a speedup of over $50\%$. We also observe that proofs with a higher original execution time tend to show more speedup. The same trends hold for miniF2F, where $114/194$ and $56/194$ proofs have a speedup over $10\%$ and $50\%$, respectively. Finally, we observe $25\%$ and $81\%$ speedups on Seed-Prover's proofs for P3 and P4 of the IMO 2025 (Sec. \ref{subsec-imo-proof-shortening}).

Upon qualitatively analyzing the proofs, we observe that the original proofs often have extraneous tactics that are eliminated by the simplified proofs. However, we also find several cases where the simplified proofs are much slower than the original proof, which usually occurs when a faster proof algorithm is replaced by a shorter but slower method (e.g. brute force with \texttt{interval\_cases}). We provide two examples of each in Appendix \ref{sec:appendix-proof-speedup}. Finally, we remark that all of our training and inference pipelines can also be applied to proof speedup as well by adjusting the reward function from proof length to proof execution time. 

\begin{figure}[htbp]
    \centering
    \begin{subfigure}{0.24\textwidth}
        \centering
        \includegraphics[width=\linewidth]{figures/speedup/scatter_putnam.pdf}
    \end{subfigure}
    \hfill
    \begin{subfigure}{0.24\textwidth}
        \centering
        \includegraphics[width=\linewidth]{figures/speedup/histogram_putnam.pdf}
    \end{subfigure}
    \begin{subfigure}{0.24\textwidth}
        \centering
        \includegraphics[width=\linewidth]{figures/speedup/scatter_minif2f.pdf}
    \end{subfigure}
    \hfill
    \begin{subfigure}{0.24\textwidth}
        \centering
        \includegraphics[width=\linewidth]{figures/speedup/histogram_minif2f.pdf}
    \end{subfigure}
    \caption{Simplified proofs are frequently faster than original proofs on miniF2F and PutnamBench.}
    \label{fig:speedup-plots}
\end{figure}

\subsubsection{Optimizing for Heartbeats instead of Proof Length} \label{subsubsec:heartbeats-optimization}

As we stated in Sec.~\ref{sec:task-metrics}, our complexity measure $\mathcal{L}$ generalizes beyond proof length. Next, we set $\mathcal{L}$ to be the number of Lean heartbeats\footnote{We use \texttt{\#count\_heartbeats} with \texttt{set\_option Elab.async false}}, a proxy of execution time that can run efficiently in parallel. With this metric, we run eight iterations of the same inference-time algorithm using ProofOptimizer-ExpIt. In Fig.~\ref{fig:speedup-plots-heartbeats} (a, b), we show analogous plots as earlier for miniF2F. Observe that this time, all the points are now on or below the $y=x$ line, eliminating the short but slow proofs we saw in Fig.~\ref{fig:speedup-plots}. Overall, we observe faster proofs, with 138/194 and 81/194 miniF2F proofs showing a speedup over 10\% and 50\%, respectively (compared to 114/194 and 56/194 before using the length metric). In Fig.~\ref{fig:speedup-plots-heartbeats} (c), we see that while the lengths of the proofs found with this metric are slightly longer than before, there is still considerable shortening. Finally, Fig.~\ref{fig:speedup-plots-heartbeats} (d) explains this by showing that proof length and number of heartbeats are generally correlated. In the future, optimizing for a combination of proof length and heartbeat count could lead to improvements in both readability and execution time. Full results can be found in Sec.~\ref{sec:appendix-heartbeats}.

% # Len: [334, 302, 163, 145, 139, 135, 129, 125, 112, 96]
% # Len Red: [0, 9.2, 71.3, 74.8, 75.8, 76.3, 76.9, 77.4, 79.0, 81.3]
% # Hb: [36290, 36236, 14505, 13636, 13327, 13173, 13023, 12786, 11938, 10382]
% # Hb Red: [0, 0.2, 43.3, 46.7, 48.2, 48.5, 48.8, 49.6, 51.5, 57.0]

% # Len: [1468, 1359, 1142, 1092, 1060, 1043, 1034, 1031, 974, 904]
% # Len Red: [0, 7.4, 32.2, 36.2, 38.7, 39.7, 40.5, 40.8, 44.0, 49.2]
% # Hb: [221378, 218872, 198727, 156647, 155133, 139967, 136239, 135990, 121899, 110820]
% # Hb Red: [0, 0.7, 18.5, 23.9, 26.9, 28.4, 29.5, 29.6, 34.0, 39.5]

\begin{figure}[htbp]
    \centering
    \begin{subfigure}{0.24\textwidth}
        \centering
        \includegraphics[width=\linewidth]{figures/speedup-heartbeats/scatter_minif2f.pdf}
    \end{subfigure}
    \hfill
    \begin{subfigure}{0.24\textwidth}
        \centering
        \includegraphics[width=\linewidth]{figures/speedup-heartbeats/histogram_minif2f.pdf}
    \end{subfigure}
    \begin{subfigure}{0.24\textwidth}
        \centering
        \includegraphics[width=\linewidth]{figures/speedup-heartbeats/putnam_min64_styled.pdf}
    \end{subfigure}
    \hfill
    \begin{subfigure}{0.24\textwidth}
        \centering
        \includegraphics[width=\linewidth]{figures/speedup-heartbeats/proof_length_vs_heartbeats.pdf}
    \end{subfigure}
    \caption{Using heartbeats instead of proof length as complexity measure}
    \label{fig:speedup-plots-heartbeats}
\end{figure}